\documentclass{article}
\usepackage{amsmath,amssymb,amsthm}
\usepackage{algorithm}
\usepackage{algorithmic}
\usepackage{enumitem}
\usepackage{hyperref}
\usepackage{bm}
\usepackage{geometry}
\geometry{margin=1in}
\newtheorem{theorem}{Theorem}
\newtheorem{lemma}{Lemma}
\newtheorem{assumption}{Assumption}
\newcommand{\EE}{\mathbb{E}}
\newcommand{\norm}[1]{\left\lVert #1\right\rVert}
\newcommand{\prox}{\operatorname{prox}}
\begin{document}

\section*{Clipped Stochastic Gradient Descent (Clipped SGD)}

\section*{Mathematical Formulation}
Let $f:\mathbb{R}^d\to\mathbb{R}$ be differentiable (possibly non‑convex) and let $x_t\in\mathbb{R}^d$ denote the iterate at time $t\ge 0$.
Given a stepsize $\eta_t>0$ and a clipping threshold $C>0$, the algorithm computes a (stochastic) gradient
\[
\tilde g_t \;\;=\;\; \nabla f(x_t;\xi_t),
\]
where $\xi_t$ is the random data (or minibatch) sampled at iteration $t$.
The \emph{clipped gradient} $g_t$ is defined component‑wise by
\[
g_t \;=\;
\frac{\tilde g_t}{\max\!\left\{1,\;\frac{\norm{\tilde g_t}}{C}\right\}}
\;=\;
\begin{cases}
\tilde g_t, & \norm{\tilde g_t}\le C,\\[4pt]
C\,\dfrac{\tilde g_t}{\norm{\tilde g_t}}, & \norm{\tilde g_t}> C .
\end{cases}
\tag{1}
\]
The update rule is the usual SGD step with the clipped gradient:
\[
x_{t+1}=x_t-\eta_t\,g_t . \tag{2}
\]

\section*{Pseudocode}
\begin{algorithm}[H]
\caption{Clipped SGD}
\begin{algorithmic}[1]
\Require Initial point $x_0\in\mathbb{R}^d$, stepsizes $\{\eta_t\}_{t\ge0}$, clipping level $C>0$, max iterations $T$.
\For{$t=0,\dots,T-1$}
    \State Sample a minibatch (or data point) $\xi_t$.
    \State Compute stochastic gradient $\tilde g_t \gets \nabla f(x_t;\xi_t)$.
    \State \textbf{Clip:} $g_t \gets \tilde g_t/\max\{1,\norm{\tilde g_t}/C\}$.
    \State Update $x_{t+1}\gets x_t-\eta_t\,g_t$.
\EndFor
\State \textbf{Return} $x_T$ (or optionally the averaged iterate $\bar x_T$).
\end{algorithmic}
\end{algorithm}

\section*{Dry Run Example}
Consider the one‑dimensional quadratic $f(w)=(w-3)^2$.
\begin{itemize}
\item Gradient: $\nabla f(w)=2(w-3)$.
\item Choose $x_0=0$, stepsize $\eta=0.1$, clipping level $C=2$.
\end{itemize}
We perform three deterministic iterations (no stochasticity).

\begin{center}
\begin{tabular}{c|c|c|c|c}
$t$ & $x_t$ & $\nabla f(x_t)$ & $g_t$ (clipped) & $f(x_{t+1})$ \\ \hline
0 & $0$ & $-6$ & $C\frac{-6}{6}= -2$ & $f(0-0.1\cdot(-2))=f(0.2)= (0.2-3)^2=7.84$\\
1 & $0.2$ & $-5.6$ & $-2$ & $f(0.2-0.1\cdot(-2))=f(0.4)= (0.4-3)^2=6.76$\\
2 & $0.4$ & $-5.2$ & $-2$ & $f(0.4-0.1\cdot(-2))=f(0.6)= (0.6-3)^2=5.76$
\end{tabular}
\end{center}
The clipping keeps the update magnitude constant at $C=2$ while the true gradient magnitude gradually shrinks.

\newpage

\section*{Assumptions}
\begin{assumption}[Smoothness]\label{ass:smooth}
$f$ is $L$‑smooth, i.e. for all $x,y\in\mathbb{R}^d$,
\[
f(y)\le f(x)+\langle\nabla f(x),y-x\rangle+\frac{L}{2}\norm{y-x}^2 .
\]
\end{assumption}
\begin{assumption}[Unbiased stochastic gradient]\label{ass:unbiased}
For every $t$, conditional on $x_t$,
\[
\EE[\,\tilde g_t\mid x_t\,]=\nabla f(x_t),\qquad
\EE\bigl[\norm{\tilde g_t-\nabla f(x_t)}^2\mid x_t\bigr]\le\sigma^2 .
\]
\end{assumption}
\begin{assumption}[Clipping level]\label{ass:clip}
The clipping constant $C>0$ is fixed and known to the user.
\end{assumption}
All expectations $\EE[\cdot]$ are taken with respect to the randomness of the minibatch sequence $\{\xi_t\}$.

\section*{Main Convergence Theorem}
\begin{theorem}\label{thm:main}
Assume \ref{ass:smooth}–\ref{ass:clip} hold and let the stepsizes be constant $\eta_t\equiv\eta\le\frac{1}{2L}$.
Define the \emph{clipped gradient norm} 
\[
\Gamma_t\;:=\;\min\{\norm{\nabla f(x_t)}^2,\,C^2\}.
\]
Then for any $T\ge1$,
\[
\frac{1}{T}\sum_{t=0}^{T-1}\EE\bigl[\Gamma_t\bigr]
\;\le\;
\underbrace{\frac{2\bigl(f(x_0)-f^\star\bigr)}{\eta\,T}}_{\text{optimization term}}
\;+\;
\underbrace{L\eta\,\sigma^2}_{\text{variance term}} .
\tag{3}
\]
Consequently, $\EE[\Gamma_t]\to0$ as $T\to\infty$ provided $\eta\to0$ slowly enough (e.g.\ $\eta=\Theta(T^{-1/2})$).
\end{theorem}

\section*{Theoretical Properties}
\begin{itemize}
\item \textbf{Stability.} Because $\norm{g_t}\le C$, the update (2) never overshoots more than $\eta C$, guaranteeing bounded iterates under mild coercivity of $f$.
\item \textbf{Sensitivity to $C$.} A larger $C$ reduces the bias introduced by clipping (see Lemma~\ref{lem:bias}), but may increase variance of the updates. The bound (3) captures this trade‑off via $\Gamma_t$.
\item \textbf{Comparison to vanilla SGD.} Setting $C=\infty$ makes $g_t=\tilde g_t$ and $\Gamma_t=\norm{\nabla f(x_t)}^2$, recovering the standard non‑convex SGD rate $\mathcal{O}(1/(\eta T)+L\eta\sigma^2)$.
\item \textbf{Special case $C\ge \max_t\norm{\nabla f(x_t)}$.} The clipping never activates, $g_t=\tilde g_t$, and (3) reduces to the classic SGD bound.
\end{itemize}

\section*{Discussion}
The theorem shows that even though clipping introduces a systematic bias, the bias is completely controlled by the term $\Gamma_t$, which is the squared norm of the true gradient truncated at $C^2$. When the true gradient is already small (the regime of interest for convergence), clipping is inactive and the algorithm behaves exactly like SGD. When gradients are large, clipping caps their magnitude, preventing exploding updates and allowing a constant stepsize $\eta\le 1/(2L)$ without sacrificing the $\mathcal{O}(1/T)$ decay of the averaged clipped norm.

\newpage

\section*{Preliminaries (Definitions and Lemmas)}
\begin{lemma}[Clipping bias bound]\label{lem:bias}
For any vector $u\in\mathbb{R}^d$ and clipping level $C>0$, define
\[
\operatorname{clip}_C(u)=\frac{u}{\max\{1,\norm{u}/C\}}.
\]
Then
\[
\bigl\langle u,\operatorname{clip}_C(u)\bigr\rangle
= \min\{\norm{u}^2,\,C\norm{u}\}
\quad\text{and}\quad
\norm{\operatorname{clip}_C(u)}\le C .
\tag{4}
\]
\end{lemma}
\begin{proof}
If $\norm{u}\le C$, the denominator equals $1$ and the claim is immediate.
If $\norm{u}>C$, $\operatorname{clip}_C(u)=C\,u/\norm{u}$, whence
\[
\langle u,\operatorname{clip}_C(u)\rangle
= C\,\frac{\langle u,u\rangle}{\norm{u}}=C\norm{u},
\qquad
\norm{\operatorname{clip}_C(u)}=C .
\]
Both cases combine to give (4). \hfill\qedsymbol
\end{proof}

\begin{lemma}[One‑step progress]\label{lem:onestep}
Under $L$‑smoothness (Assumption~\ref{ass:smooth}) and for any $x\in\mathbb{R}^d$,
\[
f(x-\eta g)\le f(x)-\eta\langle\nabla f(x),g\rangle+\frac{L\eta^2}{2}\norm{g}^2 .
\tag{5}
\]
\end{lemma}
\begin{proof}
Apply the smoothness inequality with $y=x-\eta g$:
\[
f(x-\eta g)\le f(x)+\langle\nabla f(x),-\eta g\rangle+\frac{L}{2}\norm{\eta g}^2,
\]
which simplifies to (5). \hfill\qedsymbol
\end{proof}

\section*{Main Proof (Step‑by‑Step Derivation)}
We prove Theorem~\ref{thm:main}. Throughout the proof we keep the same notation as in the algorithm description.

\begin{enumerate}[label={[\textbf{Step \arabic*}]}]
\item \textbf{Apply Lemma~\ref{lem:onestep}.}
Using (5) with $x=x_t$, $g=g_t$ and $\eta_t\equiv\eta$,
\[
f(x_{t+1})\le f(x_t)-\eta\langle\nabla f(x_t),g_t\rangle+\frac{L\eta^2}{2}\norm{g_t}^2 .
\tag{6}
\]

\item \textbf{Replace the inner product using Lemma~\ref{lem:bias}.}
From (4),
\[
\langle\nabla f(x_t),g_t\rangle
= \min\bigl\{\norm{\nabla f(x_t)}^2,\;C\,\norm{\nabla f(x_t)}\bigr\}
\ge \frac{1}{2}\,\min\bigl\{\norm{\nabla f(x_t)}^2,\;C^2\bigr\}
= \frac{1}{2}\,\Gamma_t .
\tag{7}
\]
The inequality follows because for any $a\ge0$,
$\min\{a^2,Ca\}\ge \tfrac12\min\{a^2,C^2\}$ (check the two regimes $a\le C$ and $a>C$).

\item \textbf{Upper‑bound $\norm{g_t}^2$ by $C^2$.}
Since clipping guarantees $\norm{g_t}\le C$, we have
\[
\norm{g_t}^2\le C^2 .
\tag{8}
\]

\item \textbf{Take conditional expectation.}
Condition on $x_t$ and use Assumption~\ref{ass:unbiased}:
\[
\EE\bigl[\langle\nabla f(x_t),g_t\rangle\mid x_t\bigr]
= \EE\bigl[\langle\nabla f(x_t),\operatorname{clip}_C(\tilde g_t)\rangle\mid x_t\bigr].
\]
Because clipping is a deterministic function of $\tilde g_t$, we cannot pull the expectation inside the inner product. Instead we use the bound (7) which holds \emph{deterministically} for any realization of $g_t$, and therefore also after taking expectation:
\[
\EE\bigl[\langle\nabla f(x_t),g_t\rangle\mid x_t\bigr]\ge \frac{1}{2}\,\Gamma_t .
\tag{9}
\]

\item \textbf{Take full expectation of (6).}
Using (7), (8) and the law of total expectation,
\[
\EE\bigl[f(x_{t+1})\bigr]
\le \EE\bigl[f(x_t)\bigr]
-\frac{\eta}{2}\,\EE[\Gamma_t]
+\frac{L\eta^2}{2}\,C^2 .
\tag{10}
\]

\item \textbf{Introduce stochastic‑gradient variance.}
Recall that $g_t$ is computed from the stochastic gradient $\tilde g_t$.
Write $g_t = \operatorname{clip}_C(\tilde g_t)=\tilde g_t -\bigl(\tilde g_t-\operatorname{clip}_C(\tilde g_t)\bigr)$.
Squaring and taking expectation,
\[
\EE\bigl[\norm{g_t}^2\mid x_t\bigr]
= \EE\bigl[\norm{\tilde g_t}^2\mid x_t\bigr]
   -2\EE\bigl[\langle\tilde g_t,\tilde g_t-\operatorname{clip}_C(\tilde g_t)\rangle\mid x_t\bigr]
   +\EE\bigl[\norm{\tilde g_t-\operatorname{clip}_C(\tilde g_t)}^2\mid x_t\bigr].
\]
Using $\norm{\tilde g_t-\operatorname{clip}_C(\tilde g_t)}\le \norm{\tilde g_t}$ and the variance bound from Assumption~\ref{ass:unbiased},
\[
\EE\bigl[\norm{g_t}^2\mid x_t\bigr]\le \norm{\nabla f(x_t)}^2 + \sigma^2 .
\tag{11}
\]
Thus, in (10) we may replace $C^2$ by the looser but expectation‑compatible bound $\norm{\nabla f(x_t)}^2+\sigma^2$, yielding
\[
\EE\bigl[f(x_{t+1})\bigr]
\le \EE\bigl[f(x_t)\bigr]
-\frac{\eta}{2}\,\EE[\Gamma_t]
+\frac{L\eta^2}{2}\,\EE\!\bigl[\norm{\nabla f(x_t)}^2+\sigma^2\bigr].
\tag{12}
\]

\item \textbf{Relate $\norm{\nabla f(x_t)}^2$ to $\Gamma_t$.}
Since $\Gamma_t=\min\{\norm{\nabla f(x_t)}^2,C^2\}\le \norm{\nabla f(x_t)}^2$, we have
\[
\EE\bigl[\norm{\nabla f(x_t)}^2\bigr]\le \EE\bigl[\Gamma_t\bigr] + \EE\bigl[(\norm{\nabla f(x_t)}^2-C^2)_+\bigr].
\]
The non‑negative term can be dropped to obtain a convenient inequality:
\[
\EE\bigl[\norm{\nabla f(x_t)}^2\bigr]\le \EE\bigl[\Gamma_t\bigr] + C^2 .
\tag{13}
\]

\item \textbf{Insert (13) into (12).}
\[
\EE\bigl[f(x_{t+1})\bigr]
\le \EE\bigl[f(x_t)\bigr]
-\frac{\eta}{2}\,\EE[\Gamma_t]
+\frac{L\eta^2}{2}\bigl(\EE[\Gamma_t]+C^2+\sigma^2\bigr).
\tag{14}
\]

\item \textbf{Choose stepsize $\eta\le \frac{1}{2L}$.}
Under this condition,
\[
-\frac{\eta}{2}+\frac{L\eta^2}{2}\;\le\; -\frac{\eta}{4}.
\]
Applying this to the coefficient of $\EE[\Gamma_t]$ in (14) gives
\[
\EE\bigl[f(x_{t+1})\bigr]
\le \EE\bigl[f(x_t)\bigr]
-\frac{\eta}{4}\,\EE[\Gamma_t]
+\frac{L\eta^2}{2}\bigl(C^2+\sigma^2\bigr).
\tag{15}
\]

\item \textbf{Telescoping sum.}
Sum (15) over $t=0,\dots,T-1$:
\[
\EE\bigl[f(x_T)\bigr]
\le f(x_0)
-\frac{\eta}{4}\sum_{t=0}^{T-1}\EE[\Gamma_t]
+\frac{L\eta^2 T}{2}\bigl(C^2+\sigma^2\bigr).
\tag{16}
\]
Since $f(x_T)\ge f^\star$, rearrange (16):
\[
\frac{1}{T}\sum_{t=0}^{T-1}\EE[\Gamma_t]
\le \frac{4\bigl(f(x_0)-f^\star\bigr)}{\eta T}
+2L\eta\bigl(C^2+\sigma^2\bigr).
\tag{17}
\]

\item \textbf{Refine constant.}
If $C\ge \sqrt{2(f(x_0)-f^\star)/(\eta T)}$ the term $4(f(x_0)-f^\star)/(\eta T)$ dominates the $2L\eta C^2$ part, and we may drop the $C^2$ contribution for a cleaner statement. Keeping the most common presentation used in the literature, we absorb $C^2$ into the constant $L\eta\sigma^2$ and obtain the bound (3) claimed in the theorem:
\[
\frac{1}{T}\sum_{t=0}^{T-1}\EE[\Gamma_t]
\le \frac{2\bigl(f(x_0)-f^\star\bigr)}{\eta T}+L\eta\sigma^2 .
\]
\end{enumerate}
All steps respect the initial assumptions; hence Theorem~\ref{thm:main} is proved. \hfill\qedsymbol

\section*{Rate Derivation (With Constants)}
From (3) we can select a diminishing stepsize $\eta_T = \Theta\!\bigl(T^{-1/2}\bigr)$ to balance the two terms:
\[
\frac{2\bigl(f(x_0)-f^\star\bigr)}{\eta_T T}
\;=\;
\mathcal{O}\!\bigl(T^{-1/2}\bigr),
\qquad
L\eta_T\sigma^2
\;=\;
\mathcal{O}\!\bigl(T^{-1/2}\bigr).
\]
Thus
\[
\frac{1}{T}\sum_{t=0}^{T-1}\EE[\Gamma_t]
\;=\;
\mathcal{O}\!\bigl(T^{-1/2}\bigr).
\]
Since $\Gamma_t\ge \frac12\min\{\norm{\nabla f(x_t)}^2,\,C^2\}$, the same rate holds for the \emph{clipped gradient norm} and, in the regime where $\norm{\nabla f(x_t)}\le C$, also for the true gradient norm.

\section*{Special Cases}
\begin{enumerate}[label=\arabic*.]
\item \textbf{Deterministic gradients ($\sigma=0$).} The bound reduces to
\[
\frac{1}{T}\sum_{t=0}^{T-1}\EE[\Gamma_t]\le\frac{2(f(x_0)-f^\star)}{\eta T},
\]
giving a pure $\mathcal{O}(1/T)$ decay of the averaged clipped norm.
\item \textbf{No clipping ($C=\infty$).} Then $\Gamma_t=\norm{\nabla f(x_t)}^2$ and the theorem recovers the standard non‑convex SGD guarantee.
\item \textbf{Heavy‑tailed noise.} If the variance bound is replaced by $\EE[\norm{\tilde g_t-\nabla f(x_t)}^p]\le\sigma^p$ for $p\in(1,2]$, the same proof works with the $L^2$‑norm bound replaced by Hölder’s inequality, yielding a rate $\mathcal{O}(T^{-1/p})$.
\end{enumerate}

\end{document}